\documentclass[10pt,a4paper]{amsart}
\usepackage[margin=19mm]{geometry}
\usepackage{amsmath,amssymb,mathtools}
\usepackage[scr=rsfs,cal=euler]{mathalfa}
\usepackage{xfrac}
\usepackage[unicode,hidelinks,bookmarks=false]{hyperref}
\newcommand{\ie}{i.e.\ }
\newcommand{\cf}{cf.\ }
\newcommand{\resp}{respectively\ }
\newcommand{\Subsection}{\subsection}
\providecommand{\nobreakdash}{\nobreak\hbox{-}}
\setlength{\emergencystretch}{2em}
\multlinegap0pt
\usepackage{amssymb}
\usepackage{mathtools}
\usepackage{tikz}
\usepackage{algorithm}\usepackage{algpseudocode}
\usepackage[normalem]{ulem}
\usepackage{cancel}
\newtheorem{lemma}{Lemma}
\title{\bf Vertex-critical $(P_5,\text{chair})$-free and $(P_5,\text{cricket})$-free graphs}
\author{Jorik Jooken}
\date{Source: arXiv:2605.28537v1 (CC BY 4.0)}
\begin{document}
\maketitle

\noindent\emph{Source and licence.} \bf Vertex-critical $(P_5,\text{chair})$-free and $(P_5,\text{cricket})$-free graphs. Version arXiv:2605.28537v1 (CC BY 4.0), distributed by arXiv under the Creative Commons Attribution 4.0 International licence, http://creativecommons.org/licenses/by/4.0/, as recorded in the arXiv licence metadata for this exact version and shown on the abstract page https://arxiv.org/abs/2605.28537v1. Full licence notice: \texttt{licenses/arxiv-2605.28537-CC-BY-4.0.txt}.\par\smallskip

\noindent\textit{Editorial scope.} Counted unit: the article's lemma lem:RamseyCricket (source0.tex lines 421--429). The article's complete original proof of it is delivered with it (source0.tex lines 430--440, one \texttt{proof} environment of 11 lines). No mathematics is written, repaired or re-derived: the statement and the proof are the article's own lines, in the article's own notation, and no retained line is substituted or re-flowed.  No further source-local context is needed: every cross-reference of every retained line resolves inside the two delivered spans.  Nothing was omitted from any retained span, so the package carries no omission record.  No retained line contains a citation, so the wrapper carries no bibliography and no bibliography entry.  No retained line contains a figure environment, an image-inclusion command or drawing code, so no image is delivered and the package declares no asset dependency.  Editorial formatting only, in addition: an AMS \texttt{amsart} preamble that restates the article's own declarations and macros used by the retained lines, namely the environments \texttt{lemma} on separate counters, together with the standard packages those lines need.  The retained environments print in this document as Lemma 1 in order of appearance, whereas the article prints the same items with the numbering of its own full document.  The complete native source of the article as distributed in the arXiv e-print for this exact version is bundled unchanged as \texttt{sources/201544/source0.tex}.
\begin{lemma}
\label{lem:RamseyCricket}
Let $q \geq 1$ be an integer, let $X$ be a finite set and let $Y$ be a graph with $\chi(Y) \leq q$. For every vertex $y \in Y$, a set $I(y) \subseteq X$ is given such that the following conditions hold:
\begin{itemize}
    \item If $y,y' \in Y$ are nonadjacent, then $I(y)$ and $I(y')$ are comparable.
    \item If $y,y' \in Y$ are adjacent, then $|I(y) \setminus I(y')| \leq 1$ and $|I(y') \setminus I(y)| \leq 1$.
\end{itemize}
For every $x \in X$, define $S_x := \{y \in Y~|~x \in I(y)\}$. If $\{S_x~|~x \in X\}$ forms an antichain, then $|X| \leq R_{q^2}(4)-1$.
\end{lemma}

\begin{proof}
Consider a proper $q$-colouring of $Y$ and denote the corresponding colour classes by $Y_1, Y_2, \ldots, Y_q$. Because of the first condition, we may assume that for each $1 \leq i \leq q$ the vertices in $Y_i$ are ordered such that if $y$ comes before $y'$ in this ordering, then $I(y) \subseteq I(y')$. In what follows, we will treat $Y_i$ with respect to this ordering.

For every $x \in X$ and every integer $1 \leq i \leq q$, we define $S_x^i := S_x \cap Y_i$. Note that $S_x^i$ can be obtained by removing zero or more vertices from the beginning of $Y_i$, because if $y,y' \in Y_i$ and $I(y)\subseteq I(y')$ and $x \in I(y)$, then we also have $x \in I(y')$. We shall abbreviate this observation as $(\dagger)$.

Suppose for the sake of obtaining a contradiction that $|X| \geq R_{q^2}(4)$. Let us order the elements of $X$ in an arbitrary way. If $x$ and $x'$ are two distinct elements from $X$, we write $x<x'$ if $x$ comes before $x'$ in this ordering. Let $G$ be a complete graph with vertex set $X$. Let $x, x' \in X$ be two distinct elements such that $x<x'$. Recall that $S_x$ and $S_{x'}$ are incomparable, because $\{S_x~|~x \in X\}$ forms an antichain. Therefore, there exists an integer $i$ such that $S_x^i \setminus S_{x'}^i \neq \emptyset$ and an integer $j$ such that $S_{x'}^j \setminus S_x^j \neq \emptyset$. We colour the edge $xx' \in E(G)$ with colour $(i,j)$. Note that there are at most $q^2$ different colours.

Since $|X| \geq R_{q^2}(4)$, there are four vertices $x_1, x_2, x_3, x_4 \in V(G)$ such that each edge between these four vertices is coloured with the same colour. Let us call this colour $(i,j)$. By the definition of the colours, we have for every $1 \leq s < t \leq 4$ that $S_{x_s}^i \setminus S_{x_t}^i \neq \emptyset$ and $S_{x_t}^j \setminus S_{x_s}^j \neq \emptyset$. By observation $(\dagger)$, we have $S_{x_t}^i \subset S_{x_s}^i$ and also $S_{x_s}^j \subset S_{x_t}^j$. Hence, we have $S_{x_4}^i \subset S_{x_3}^i \subset S_{x_2}^i \subset S_{x_1}^i$ and $S_{x_1}^j \subset S_{x_2}^j \subset S_{x_3}^j \subset S_{x_4}^j$.

Consider a vertex $y \in S_{x_2}^i \setminus S_{x_3}^i$ and a vertex $y' \in S_{x_3}^j \setminus S_{x_2}^j$. Note that we have $x_1, x_2 \in I(y)$, but $x_3, x_4 \notin I(y)$ and $x_3, x_4 \in I(y')$, but $x_1, x_2 \notin I(y')$. Therefore we have $|I(y) \setminus I(y')| \geq 2$ and $|I(y') \setminus I(y)| \geq 2$. This contradicts the two conditions and proves the lemma.
\end{proof}

\end{document}
